

\documentclass[11pt]{article}

\usepackage[utf8]{inputenc}
\usepackage[T1]{fontenc}
\usepackage[margin=1in]{geometry}
\usepackage{hyperref}
\usepackage{url}
\usepackage{booktabs}
\usepackage{amsmath}
\usepackage{amssymb}
\usepackage{xcolor}
\usepackage{enumitem}
\usepackage{microtype}
\usepackage{graphicx}
\usepackage[numbers,sort&compress]{natbib}

\newcommand{\zvf}{\textsc{Zvf}}
\newcommand{\gu}{\textsc{Gu}}
\newcommand{\minreport}{\textsc{Min-Report-RL}}
\newcommand{\grpo}{\textsc{GRPO}}
\newcommand{\audit}{\textsc{GRPO-Survival-Audit}}
\newcommand{\grpoReg}{\textsc{GRPO-Registry}}

\title{\audit{}: A Single-Stack Survival Protocol for the \grpo{} Family\\
\large Holding the Stack Fixed So Algorithmic Gains Can Be Seen}

\author{Arvind C R\thanks{Pillar~3 of the ZVF Program.}\\
PES University\\\texttt{arvindcr4@gmail.com}}

\date{\today}

\begin{document}
\maketitle

\begin{abstract}
The \grpo{}-family literature reports head-to-head comparisons in which each
variant runs in its own trainer, sampler, and evaluation harness. Because the
``stack''---loss form, reference/KL handling, sampler precision, group-size
schedule, \zvf/\gu{} trajectory, held-out split, and parser---is a documented
flip lever \cite{minreportrl2026}, such comparisons confound algorithmic
innovation with implementation differences. This paper proposes \audit{}, a
single-stack reproducibility-audit protocol in which DAPO, GSPO, Dr.\grpo{},
and AERO are re-implemented as declared overrides on one shared language-model
trainer. M-\grpo{} is assigned to a separate agentic stratum because
hierarchical credit assignment changes the environment and trajectory schema.
We specify the shared
stack, the pre-registration rules, the survival-verdict logic, the statistical
protocol, and the role of the companion \texttt{grpo-stackdiff} tool in
verifying that arms are stack-matched. We also report a small open-source pilot
on Qwen2.5-0.5B-Instruct (two seeds, T4 GPU) that demonstrates the protocol's
feasibility and produces concrete \zvf{} and held-out-delta observations. The
full-scale frozen audit is complete: all 40 arm--seed units independently pass
the local, W\&B, private-Hub checkpoint, stack-fingerprint,
treatment-fingerprint, and 500-row held-out gates. Against the shared GRPO
reference, DAPO's controlled delta is $+0.0015$ (paired 95\% CI
$[-0.0040,0.00725]$) and receives the preregistered
\texttt{DISAPPEARS} verdict. GSPO, Dr.\grpo{}, and AERO receive
\texttt{INCONCLUSIVE} verdicts; none of the five arms collapsed. These results
are a stack-controlled survival audit, not a universal capability leaderboard.
\end{abstract}

\section{Introduction: The Need for a Survival Audit}


A typical GRPO-family paper compares a new variant against ``GRPO'' by running
each in a different codebase. The new variant may change the clip rule, the
advantage normalization, or the group-size schedule; the baseline may use a
different token mask, KL placement, sampler precision, or reward parser
\cite{minreportrl2026, zvfaudit2026}. The result is a comparison of two stacks
$S_A$ and $S_B$ that happen to be tagged ``A'' and ``B''. The algorithmic delta
and the stack delta are inseparable.

The \minreport{} position paper \cite{minreportrl2026} proposes an eight-item
standard---seven run-manifest fields plus held-out pass@$k$ reporting---so that future papers expose the
stack. The \grpoReg{} catalog \cite{grpoRegistry2026} records each variant's
self-declared delta relative to a baseline. This paper closes the loop: it
specifies a controlled, single-stack re-implementation protocol that asks,
for each variant, \emph{what fraction of its claimed gain survives when the
stack is held fixed?} We call the answer the variant's \emph{survival}.

\paragraph{What survival is not.} A low survival verdict does not mean a
variant is ``fake.'' It means the variant's published gain was partly or wholly
a stack effect under the audited stack. A high survival verdict is evidence
that the algorithmic idea is robust. Both outcomes are useful.

\paragraph{A live specimen.} While assembling this audit's own pilot, our
trainer shipped exactly the failure the protocol targets: a
\texttt{--loss drgrpo} flag whose commit message claimed it toggled the
advantage normalization, but which was never wired into the loss --- so a
six-arm GRPO-vs-Dr.GRPO panel silently compared GRPO with GRPO, and its
committed ``finding'' stood until the runner was read line by line. No
output-level check caught it: rewards, lengths, and ZVF traces all looked
plausible. The panel was invalidated, the artifacts preserved under
explicit \texttt{.invalid\_actually\_grpo.json} names, and the corrected
rerun (2026-07-11) reversed one of the two claims. Stack identity cannot
be certified from results; it must be certified from the stack
(\S\ref{sec:stackdiff}).

\paragraph{Contributions.}
\begin{enumerate}[leftmargin=1.6em, itemsep=2pt]
  \item A concrete single-stack audit protocol with shared-stack constraints,
    arm definitions, and pre-registration rules (\S\ref{sec:design}).
  \item A survival-verdict logic and statistical protocol that treat seed as
    the independent unit (\S\ref{sec:verdict}--\ref{sec:stats}).
  \item Integration with \texttt{grpo-stackdiff} to verify stack identity
    across arms before any result is trusted (\S\ref{sec:stackdiff}).
  \item A small open-source pilot (\S\ref{sec:pilot}) that exercises the
    protocol on four arms and shows that the required telemetry is obtainable
    from a standard GRPO trainer.
  \item A fail-closed execution snapshot: all eight independently reconciled
    GRPO and GSPO seed records, DAPO seeds 11, 23, 53, 71, 89, 107, and 131,
    Dr.\grpo{} seeds 11, 37, 53, 71, 89, 107, and 131, AERO seeds 11, 23, 37,
    53, 71, 89, and 107, three explicitly missing units, and no
    premature arm-level verdict.
\end{enumerate}

\section{Audit Design}
\label{sec:design}


\paragraph{Shared stack.} Every arm uses the same:
\begin{itemize}[leftmargin=1.4em, itemsep=2pt]
  \item base checkpoint and tokenizer revision,
  \item LoRA target modules, rank, alpha, and dropout (if LoRA is used),
  \item optimizer, learning-rate schedule, and gradient clipping,
  \item rollout engine (vLLM), decoding parameters, sampler/trainer precision,
  \item reference-policy and KL handling (placement, coefficient, schedule,
    estimator),
  \item token mask and advantage normalization (unless the variant explicitly
    changes one of these),
  \item reward parser and held-out evaluation harness,
  \item prompt distribution and train/test split.
\end{itemize}
The only permitted differences are the variant's declared algorithmic deltas.

\paragraph{Arms.} The primary arms are:
\begin{itemize}[leftmargin=1.4em, itemsep=2pt]
  \item \textbf{GRPO (baseline):} vanilla per-token GRPO with the shared stack.
  \item \textbf{DAPO:} asymmetric clip-higher + dynamic sampling
    \cite{yu2025dapo}.
  \item \textbf{GSPO:} sequence-level importance-sampling ratio
    \cite{qwen2025gspo}.
  \item \textbf{Dr.\grpo{}:} removal of within-group standard-deviation and
    length-dependent normalization
    \cite{tong2025drgrpo}.
  \item \textbf{AERO:} adaptive rollout allocation, selective rejection, and
    posterior-guided avoidance of zero-advantage groups \cite{zhang2026aero}.
\end{itemize}
Optional secondary arms (CPPO \cite{lin2025cppo}, NGRPO
\cite{nan2025ngrpo}, Scaf-\grpo{} \cite{zhang2025scaffgrpo}) are added if compute
permits; each would be implemented as a declared hook on the same shared trainer.
M-\grpo{} \cite{hong2025mgrpo} is audited separately on an agentic task with a
hierarchical trajectory schema; pooling it with arithmetic arms would violate
the stack-lock requirement. Its separate frozen contract is
\path{zvf-program/audit/preregistration_mgrpo.json}.

\paragraph{Pre-registration.} The full audit is locked to Qwen3-8B, GSM8K,
30 training steps, batch size 4, 1,024-token completions, a fixed held-out set
of 500 prompts, and eight paired seeds
$\{11,23,37,53,71,89,107,131\}$. The checkpoint is the fixed final step, not
the best training-reward checkpoint. The primary estimand is the seed-paired
held-out delta relative to GRPO; training reward, \zvf/\gu{}, realized rollouts,
and wall-clock time are secondary. The machine-readable record is
\texttt{audit/preregistration.json}; any change after a scored run requires a
new version and is not pooled with this study. Pre-registration prevents the
selection-by-reward bias that the audit is designed to diagnose
\cite{zvfaudit2026}.

\paragraph{Hook interface.} A variant is added by implementing a small hook
object with the following interface:
\begin{verbatim}
class VariantHook:
    def on_group_rewards(self, rewards: Tensor) -> Tensor:
        # optionally transform per-prompt group rewards
        return rewards
    def on_advantage(self, advantage: Tensor) -> Tensor:
        # optionally transform advantages
        return advantage
    def on_loss(self, loss_terms: Tensor, mask: Tensor) -> Tensor:
        # optionally modify loss terms
        return loss_terms
    def group_size_rule(self, step: int, zvf: float) -> int:
        # return G for this step; default is fixed G
        return self.G
\end{verbatim}
All hooks are loaded into the same trainer binary, so infrastructure differences
are impossible.

\section{Survival Verdicts}
\label{sec:verdict}


For each variant we report two numbers when they are metric-compatible:
\begin{itemize}[leftmargin=1.4em, itemsep=2pt]
  \item \textbf{Published $\Delta$:} the gain claimed in the original paper,
    used quantitatively only when model, task, split, and metric match the
    controlled audit. Otherwise it is contextual metadata.
  \item \textbf{Controlled $\Delta$:} the gain in the shared stack vs.\ the
    shared GRPO baseline, with a 95\% confidence interval across seeds.
\end{itemize}
The audit verdict is then:
\begin{center}
\small
\begin{tabular}{@{}lp{0.72\linewidth}@{}}
\toprule
\textbf{Verdict} & \textbf{Rule}\\
\midrule
Retains & Controlled CI excludes zero and its lower bound is at least $50\%$
  of a metric-compatible published $\Delta$.\\
Survives & Controlled CI excludes zero, but the published effect is not
  metric-compatible or the half-gain rule fails.\\
Disappears & The 90\% CI lies inside a preregistered equivalence margin.\\
Reverses & Controlled 95\% CI is entirely negative.\\
Inconclusive & Every other outcome, including inadequate power.\\
\bottomrule
\end{tabular}
\end{center}
The $50\%$ retention threshold is a preregistered convention, not a comparison
between incommensurate benchmark deltas. The equivalence margin is 0.01 in
absolute held-out success. Sensitivity analyses at 0.005 and 0.02 are secondary.

\section{Statistical Protocol}
\label{sec:stats}


The independent unit is a seed, not a training step. Training rewards are
autocorrelated within a trace, so inferential claims are made only across the
seed sweep. For each metric we report the mean and standard error across seeds,
with bootstrap 95\% confidence intervals. We control the false-discovery rate
across the pre-registered survival tests using Benjamini--Hochberg. Single-seed
arms are reported as descriptive only.

The chosen seed count is $S=8$ per core arm. The achieved minimum detectable
effect is computed from the seed-level standard deviation before verdicts are
assigned. If the MDE exceeds the equivalence margin $m$, the audit cannot claim
disappearance; it reports \texttt{INCONCLUSIVE}.

\section{Using \texttt{grpo-stackdiff} to Verify Stack Identity}
\label{sec:stackdiff}


Before interpreting any survival result, we verify that the arms are
stack-matched. The implemented registry CLI command
\path{platform_hybrid/registry/query.py stackdiff} \cite{minreportrl2026}
takes two \minreport{} registry records and reports
whether any unapproved lever differs. The audit proceeds only if every pairwise
comparison returns \texttt{STACK\_MATCHED} (or \texttt{STACK\_MATERIAL} for
levers the protocol explicitly tolerates).

If a hook inadvertently changes a non-target lever---for example, a dynamic
sampling rule that also changes the effective tokenizer chat template because
it filters on string prefixes---\texttt{grpo-stackdiff} flags it as a
DISTRIBUTIONAL or SEMANTIC\_OBJECTIVE diff. The hook must then be rewritten so
that only the declared delta remains. This check is the difference between a
``single-trainer'' audit and a genuinely ``single-stack'' audit.

\section{Locked Outputs and Execution Status}


The locked output schema contains one row per arm and seed plus two generated
tables: (i) controlled held-out delta, paired 95\% CI, achieved MDE, and verdict;
and (ii) last-10 reward, late \zvf/\gu{}, collapse rate, realized rollout count,
group schedule, and wall-clock time. The aggregator refuses to emit a verdict
unless all eight core seeds are present, the manifest stack lock passes, and
every metric has the preregistered held-out size. As of 2026-07-19, the locked
directory contains all eight independently accepted records for each of GRPO,
DAPO, GSPO, Dr.\grpo{}, and AERO
each with 500 held-out examples. The GRPO held-out exact matches are
respectively 0.646, 0.636, 0.634, 0.634, 0.634, 0.620, 0.634, and 0.622; DAPO
seeds 11, 23, 37, 53, 71, 89, 107, and 131 achieve 0.648, 0.624, 0.628, 0.636, 0.634,
0.636, 0.634, and 0.632; GSPO seeds 11, 23, 37, 53, 71, 89, 107, and 131 achieve 0.640,
0.638, 0.642, 0.632, 0.638, 0.646, 0.640, and 0.628; and Dr.\grpo{} seeds 11, 23, 37, 53, 71, 89, 107, and
131 achieve 0.642, 0.626, 0.622, 0.636, 0.620, 0.628, 0.628, and 0.646; AERO seeds 11, 23,
37, 53, 71, 89, 107, and 131 achieve 0.630, 0.628, 0.636, 0.632, 0.632, 0.630, 0.630, and 0.640. Each accepted
record reconciles a finished W\&B run, a private Hugging Face repository with
full checkpoints at steps 5, 10, 15, 20, 25, and 30, the frozen stack and
treatment fingerprints, and the 500-row evaluation trace. The aggregator
therefore reports \texttt{COMPLETE} and emits all four preregistered paired
verdicts:

\begin{table}[t]
\centering
\caption{\textbf{Frozen E1 stack-controlled results.} Deltas and confidence
  intervals are paired by seed against GRPO; MDE is the achieved 80\%-power
  minimum detectable effect.}
\label{tab:e1-results}
\small
\begin{tabular}{@{}lrrrrl@{}}
\toprule
\textbf{Arm} & \textbf{$\Delta$} & \textbf{95\% CI low} &
\textbf{95\% CI high} & \textbf{MDE$_{80}$} & \textbf{Verdict} \\
\midrule
DAPO       & $+0.0015$  & $-0.0040$  & $+0.00725$ & $0.00859$ & \texttt{DISAPPEARS} \\
GSPO       & $+0.0055$  & $ 0.0000$  & $+0.01225$ & $0.00940$ & \texttt{INCONCLUSIVE} \\
Dr.\grpo{} & $-0.0015$  & $-0.00875$ & $+0.00750$ & $0.01252$ & \texttt{INCONCLUSIVE} \\
AERO       & $-0.00025$ & $-0.00675$ & $+0.00675$ & $0.01040$ & \texttt{INCONCLUSIVE} \\
\bottomrule
\end{tabular}
\end{table}

All observed collapse rates are zero. DAPO is the sole arm whose result meets
the preregistered disappearance rule; the other three intervals do not justify
either disappearance or survival under this eight-seed design. The descriptive
pilot below remains separately labelled infrastructure evidence. The frozen
contract and aggregator are
\path{zvf-program/audit/preregistration.json} and
\path{zvf-program/audit/aggregate_audit.py}.

For DAPO seed 131, post-acceptance review found a recovery-finalizer
bookkeeping defect that had hardcoded the secondary rollout count to 480. The
held-out trace and model artifacts were unaffected. The accepted manifest now
derives 1,728 generated rollouts from the immutable checkpoint-30 trainer state,
preserves the frozen treatment fields, and carries an explicit correction
receipt mirrored in W\&B and the private Hub repository.

For GSPO seed 71, the same older evaluation-only finalizer omitted the frozen
treatment fields from \texttt{run\_config}; its rollout count, training, and
held-out evidence were already correct. A provenance-only correction restores
those fields and records a receipt mirrored in W\&B and the private Hub
repository. The corrected local and remote manifests are byte-identical.

For AERO seed 37, the recovered 500-row evaluation was complete before an
evaluation-only finalizer rejected the checkpoint's cumulative rollout count.
The defect applied the fixed-arm minimum of 480 to AERO, whose frozen
three-initial-plus-one-rescue design generates 12--16 real rollouts per step.
The accepted 30-step telemetry value of 436 lies within the correct 360--480
bound. The correction changed only the recovery finalizer and its failure
classifier; no training or held-out example was regenerated. Independent
read-back verifies W\&B run \texttt{7547ba19}, private Hub commit
\texttt{613d3f70cc0fb89eae5775a398d6f66136d14f49}, all six checkpoint trees,
the byte-identical local/remote manifest, frozen fingerprints, and 500
contiguous hashed held-out rows with 318 correct.

\section{Open-Source Pilot}
\label{sec:pilot}


We ran a small pilot on a single T4 GPU to validate that the required telemetry
can be extracted from a standard GRPO trainer and that the protocol produces
informative stack-controlled observations. The pilot is \emph{not} the
full-scale audit; it uses a small model and only two seeds.

\paragraph{Setup.} Model: Qwen/Qwen2.5-0.5B-Instruct with full-model updates;
seeds $\{0,1\}$; ten training steps; three prompts per batch; initial group size
$G_0=4$; maximum generation length 40 tokens; and a seed-specific held-out set
of 20 generated integer-addition prompts. Arms are GRPO, Dr.\grpo{}, DAPO, and
an adaptive-$G$ GRPO baseline. The shared implementation keeps tokenizer,
sampler, reward parser, prompt generator, optimizer settings, and evaluation
procedure fixed across arms. These small synthetic held-out sets make the
pilot a telemetry smoke test, not a language-model capability benchmark.

\paragraph{Pilot results.} Table~\ref{tab:pilot} shows mean held-out delta,
mean \zvf{}, and mean rollouts per arm. Even at this scale, the stack-controlled
observations are informative: DAPO reports zero \zvf{} but spends more rollouts
(174 vs.\ 120), while adaptive-$G$ matches Dr.\grpo{}'s held-out delta with a
different \zvf{} and rollout budget.

\begin{table}[t]
\centering
\caption{\textbf{Pilot results (Qwen2.5-0.5B-Instruct, 2 seeds, T4).} Mean
  held-out $\Delta$, mean \zvf{}, and mean rollouts per arm.}
\label{tab:pilot}
\small
\begin{tabular}{@{}lccc@{}}
\toprule
\textbf{Arm} & \textbf{Held-out $\Delta$} & \textbf{Mean \zvf{}} &
\textbf{Mean rollouts}\\
\midrule
GRPO (baseline)          & $0.500$ & $0.250$ & $120$ \\
Dr.\grpo{}                & $0.575$ & $0.267$ & $120$ \\
DAPO                     & $0.550$ & $0.000$ & $174$ \\
GRPO adaptive-$G$        & $0.575$ & $0.233$ & $186$ \\
\bottomrule
\end{tabular}
\end{table}

\paragraph{Interpretation.} This two-seed pilot is descriptive and cannot rank
the arms. It shows that the same held-out delta can arrive through different
signal-budget paths. DAPO's zero recorded \zvf{} is mechanical: its resampling
loop rejects flat groups before an optimizer step, spending 174 rather than 120
rollouts. Adaptive-$G$ reaches the same descriptive delta as Dr.\grpo{} while
spending 186 rollouts. The comparison therefore validates telemetry and budget
accounting only; it is not evidence that one method learns better.

\section{Limitations and Extensions}


\paragraph{Scope.} The audit tests survival under one shared stack. A variant
that survives here may still fail under a different stack; the audit does not
claim universal robustness. It claims only that the published gain is not
purely an artifact of the original stack.

\paragraph{Faithfulness of re-implementation.} A variant may depend on an
unreported implementation detail. If a careful re-implementation cannot
reproduce the original result from the paper, that is itself a finding about
reporting quality, but it is not a finding about the method's intrinsic value.
We report both readings and publish every arm's full \minreport{} block.

\paragraph{Extensions.} The protocol generalizes to any family of methods that
are small deltas on a base loop. The same hook interface can accommodate
future variants; the same \texttt{grpo-stackdiff} check can verify that each
new hook changes only its declared levers.

\section{Conclusion}


The \grpo{} family needs a way to separate algorithmic gains from stack
effects. The \minreport{} standard tells authors what to report; the
\grpoReg{} catalog records each variant's declared delta; and the \audit{}
protocol reported here specifies how to hold the stack fixed so that those
deltas can be tested fairly. The frozen multi-seed audit is complete: all 40
units passed the evidence gates, DAPO's reported advantage disappears under
the shared stack, and GSPO, Dr.\grpo{}, and AERO remain inconclusive at eight
paired seeds. The protocol, telemetry pilot, machine-readable preregistration,
complete execution record, and \texttt{grpo-stackdiff} verification step
together
they make it operational to ask of a GRPO-family result not
``did it win?'' but ``did it survive?''

\bibliographystyle{plainnat}
\bibliography{../../platform_hybrid/paper/references}

\end{document}
